\documentclass[11pt]{article}

\usepackage[margin=1in]{geometry}
\usepackage{amsmath,amssymb,amsthm}
\usepackage{graphicx}
\usepackage{booktabs}
\usepackage{siunitx}
\usepackage[colorlinks=true,linkcolor=blue,citecolor=blue]{hyperref}

\newtheorem{lemma}{Lemma}
\newtheorem{proposition}{Proposition}
\newtheorem{corollary}{Corollary}
\theoremstyle{remark}
\newtheorem{remark}{Remark}

\newcommand{\wbar}{\bar{\omega}}
\newcommand{\pbar}{\bar{\psi}}
\newcommand{\DT}{\Delta T}
\newcommand{\DTs}{\Delta T^{\star}}
\newcommand{\taul}{\tau_{\lambda}}
\newcommand{\norm}[1]{\left\lVert #1 \right\rVert}

\title{Theoretical Guarantees for the AB2CN2 Temporal Closure\\
\large Convergence radius, stencil resolution, decorrelation, and error propagation}
\author{}
\date{}

\begin{document}
\maketitle

\begin{center}
\begin{minipage}{0.92\textwidth}
\small\itshape
Coarse-grid quasi-geostrophic (QG) turbulence, machine-learned temporal
closure ($\delta R$). All results: float64 arithmetic; dealiased
($2/3$-rule) Jacobians matching the solver right-hand side; $512^2$
(FRC-256: $256^2$) forced turbulence; trajectory marks at
$dt_{\mathrm{fine}} = 5\times10^{-3}$. Reference truth is the analytic
spectral chain rule; all errors are relative $L_2$ in the resolved band.
Members: FRC-Re25k ($\nu = 4.0\times10^{-5}$, $\beta = 1$); FRC-combo
($\nu = 4.0\times10^{-5}$, $\beta = 0.5$); FRC-kf4
($\nu = 1.02\times10^{-4}$, $\beta = 1$); FRC-256
($\nu = 1.02\times10^{-4}$, $\beta = 1$, $256^2$).
\end{minipage}
\end{center}

\section*{Summary}

The closure has a $\DT$-dependent domain of validity set by two nested
limits: an \emph{outer wall} --- the truncation series has a finite
convergence radius $\DTs$ --- and an \emph{inner wall} --- a finite-lag time
stencil stops resolving the $N$-derivatives at a smaller step, and is hit
first. Both are governed by one measurable timescale, the Taylor microscale
of the nonlinear term, $\taul = 1/\sigma$ with
$\sigma = \norm{\dot N}_{\mathrm{rms}}/\norm{N}_{\mathrm{rms}}$
(derived and verified in \S\ref{sec:decorr}). Empirically
$\DTs \approx 2\,\taul$ across all measured regimes
(\S\ref{sec:radius-micro}). Error propagation is benign across the entire
trained envelope $\DT \le 1.5\times10^{-2}$: the closure error equals the
$\ddot N$ approximation error one-to-one, un-amplified
(\S\ref{sec:amplification}).

\section{Governing equation and $N$-derivatives}\label{sec:setup}

Coarse QG vorticity:
\begin{equation}
\partial_t \wbar = L\wbar + N, \qquad
N = -J(\pbar, \wbar) + F, \qquad
\pbar = \nabla^{-2}\wbar, \qquad \dot F = 0,
\end{equation}
with the spectral symbol
$\hat L(\mathbf{k}) = -\nu|\mathbf{k}|^2 - \mu + i\beta k_x/|\mathbf{k}|^2$
and $J(a,b) = \partial_x a\,\partial_y b - \partial_y a\,\partial_x b$.

Time derivatives of the state and of the nonlinear term follow from the
chain rule and the bilinearity of $J$:
\begin{equation}\label{eq:recursion}
\wbar^{(k)} = L\,\wbar^{(k-1)} + N^{(k-1)}, \qquad
\pbar^{(k)} = \nabla^{-2}\wbar^{(k)},
\end{equation}
\begin{equation}\label{eq:binomial}
N^{(m)} = -\sum_{j=0}^{m}\binom{m}{j}\,
J\!\bigl(\pbar^{(m-j)},\,\wbar^{(j)}\bigr), \qquad m \ge 1,
\end{equation}
where the static forcing drops out of every $N^{(m)}$, $m\ge1$. Every
quantity in this document is built from these exact relations.

\section{The truncation series}\label{sec:series}

The defect of one AB2CN2 step of size $\DT$ against the true flow is a power
series in $\DT$:
\begin{equation}\label{eq:defect}
\delta(\DT) = \sum_{p\ge3} c_p\,\DT^{\,p}, \qquad c_p = \frac{R_p}{D_p},
\end{equation}
\begin{align}
R_3 &= L^3\wbar + L^2 N + L\dot N - 5\ddot N, & D_3 &= 12,\\
R_4 &= 2L^4\wbar + 2L^3 N + 2L^2\dot N - 4L\ddot N + \dddot N, & D_4 &= 24,
\end{align}
with $D_5 = 240$, $D_6 = 1440$. Each coefficient is dominated by its highest
$N$-derivative $N^{(p-1)}$. The weight $-5$ on $\ddot N$ in $R_3$ is the AB2
stencil's structural blind spot to $\ddot N$; it makes the $\ddot N$ term the
dominant, irreducible part of the closure (quantified in
\S\ref{sec:amplification}).

\section{Outer wall: the convergence radius $\DTs$}\label{sec:outer}

\begin{lemma}[Cauchy--Hadamard]\label{lem:ch}
The series $\delta = \sum_p c_p \DT^{\,p}$ converges for $\DT < \DTs$ with
\begin{equation}
\frac{1}{\DTs} = \limsup_{p\to\infty}\,\norm{c_p}^{1/p}.
\end{equation}
\end{lemma}

The successive ratio $\norm{c_p}/\norm{c_{p+1}}$ only \emph{brackets} $\DTs$
(the $c_p$ are not geometric); the headline number is the root-test median
over $p = 3,\dots,6$.

\paragraph{Measurement} ($p = 3,\dots,6$, 32 anchors per member):

\begin{center}\small
\begin{tabular}{lcccccc}
\toprule
member & $\nu$ & $\beta$ &
$\norm{N^{(m+1)}}/\norm{N^{(m)}}$, $m{=}0{\to}1,\dots,4{\to}5$ &
$\DTs$ (root) & ratio bracket\\
\midrule
Re25k & $4.0{\times}10^{-5}$ & 1.0 & 34, 58, 72, 84, 93 & $\mathbf{0.066}$ & $[0.017,\,0.136]$\\
combo & $4.0{\times}10^{-5}$ & 0.5 & 15, 27, 37, 45, 51 & $\mathbf{0.139}$ & $[0.033,\,0.270]$\\
kf4   & $1.0{\times}10^{-4}$ & 1.0 & 10, 19, 28, 33, 37 & $\mathbf{0.199}$ & $[0.043,\,0.366]$\\
\bottomrule
\end{tabular}
\end{center}

\begin{proposition}\label{prop:radius}
$\DTs$ is finite and set by the growth of the $N$-derivatives. For
$\DT > \DTs$ no finite-order analytic closure converges.
\end{proposition}

\paragraph{$\beta$-dependence at fixed $Re$.}
Re25k and combo share $\nu$ and differ only in $\beta$ ($1.0$ vs $0.5$):
halving $\beta$ roughly doubles $\DTs$ ($0.066 \to 0.139$). The radius
depends on $(Re,\beta)$ jointly through the flow's rate content, made
precise in \S\ref{sec:decorr}--\ref{sec:radius-micro}. Note also that the
per-order growth ratios in the table \emph{increase with} $m$
($34 \to 93$ for Re25k): higher derivatives grow faster than any single
rate. \S\ref{sec:decorr} derives why.

\section{Inner wall: the finite-lag time stencil}\label{sec:inner}

The $N$-derivatives entering the closure are estimated from $k$ stored marks
at lags $\{0,\, \DT,\, \dots,\, (k{-}1)\DT\}$\\linebreak[1] (a backward finite difference). Two
opposing effects of stencil depth:

\begin{center}\small
\begin{tabular}{ll}
\toprule
effect & direction\\
\midrule
higher approximation order, $O(\DT^{\,k-m})$ for the $m$-th derivative & deeper helps\\
longer reach back in time, $(k{-}1)\DT$ & deeper hurts once old marks are stale\\
\bottomrule
\end{tabular}
\end{center}

\paragraph{Measurement.}
Relative error of $\ddot N$ from the depth-4 vs depth-7 stencil, exact
spectral spatial operators, three regimes $\times$ three steps:

\begin{center}
\begin{tabular}{lcccl}
\toprule
member & $\DT$ & $\ddot N$ error, $k{=}4$ & $\ddot N$ error, $k{=}7$ & deeper stencil\ldots\\
\midrule
Re25k & 0.005 & 0.149 & $\mathbf{0.020}$ & helps $7.6\times$\\
Re25k & 0.010 & 0.500 & $\mathbf{0.411}$ & helps, marginal\\
Re25k & 0.015 & $\mathbf{0.987}$ & 1.510 & \textbf{hurts}\\
combo & 0.005 & 0.033 & $\mathbf{0.001}$ & helps $25\times$\\
combo & 0.010 & 0.126 & $\mathbf{0.014}$ & helps $9\times$\\
combo & 0.015 & 0.273 & $\mathbf{0.092}$ & helps $3\times$\\
kf4   & 0.005 & 0.021 & $\mathbf{0.003}$ & helps $7.6\times$\\
kf4   & 0.010 & 0.084 & $\mathbf{0.004}$ & helps $19\times$\\
kf4   & 0.015 & 0.183 & $\mathbf{0.031}$ & helps $6\times$\\
\bottomrule
\end{tabular}
\end{center}

The crossover is directly observed: for the narrow-radius member (Re25k,
$\DTs = 0.066$) the depth-7 stencil helps at $\DT = 5\times10^{-3}$, barely
helps at $10^{-2}$, and \emph{hurts} at $1.5\times10^{-2}$ --- where its span
$6\DT = 0.09$ exceeds $\DTs$. For the wide-radius members (combo, kf4)
depth-7 wins at every measured step. The inner wall is where a stencil's
span becomes comparable to the member's $\DTs$; its precise location is
regime-dependent and, for combo/kf4, lies beyond the measured range.

\begin{proposition}\label{prop:inner}
Deeper stencils lower the $N$-derivative error while the stencil span stays
well inside $\DTs$, and raise it once the span exceeds $\DTs$. Observed at
Re25k: help at $5\times10^{-3}$ (span $0.03 \ll 0.066$), harm at
$1.5\times10^{-2}$ (span $0.09 > 0.066$).
\end{proposition}

\section{The correlation function and the Taylor microscale}\label{sec:decorr}

The two walls above invoke ``staleness'' of old marks. This section makes
that precise, from first principles, and verifies it.

\paragraph{Setup.}
Along a developed-flow trajectory, view the nonlinear term as a process
$N(t)$ with the domain inner product $\langle a,b\rangle = \int_\Omega a\,b\,dx$.
Assume statistical stationarity (S): $\mathbb{E}\norm{N(t)}^2$ constant over
the analysis window (spin-up windows are excluded). Define the
autocorrelation
\begin{equation}
\rho(\tau) = \frac{\mathbb{E}\,\langle N(t),\,N(t+\tau)\rangle}
                  {\mathbb{E}\,\norm{N(t)}^2}.
\end{equation}

\paragraph{Derivation.}
Taylor-expand $N(t+\tau)$ inside the correlation:
\begin{equation}
\mathbb{E}\langle N, N(t+\tau)\rangle
= \mathbb{E}\norm{N}^2
+ \tau\,\mathbb{E}\langle N,\dot N\rangle
+ \frac{\tau^2}{2}\,\mathbb{E}\langle N,\ddot N\rangle
+ O(\tau^3).
\end{equation}
The first-order term vanishes:
$\langle N,\dot N\rangle = \tfrac12\,\frac{d}{dt}\norm{N}^2$, whose
stationary average is zero. Differentiating
$\mathbb{E}\langle N,\dot N\rangle = 0$ once more in $t$ gives
$\mathbb{E}\langle N,\ddot N\rangle = -\,\mathbb{E}\norm{\dot N}^2$. Hence
\begin{equation}\label{eq:identity}
\boxed{\;
\rho(\tau) = 1 - \frac{\tau^2}{2}\,\sigma^2 + O(\tau^4),
\qquad
\sigma^2 = \frac{\mathbb{E}\norm{\dot N}^2}{\mathbb{E}\norm{N}^2},
\qquad
\taul = \frac{1}{\sigma},
\;}
\end{equation}
an exact identity under (S) --- no Gaussianity, no spectral assumption.
$\taul$ is the \emph{Taylor microscale} of the $N$-process. The same
derivation holds verbatim for the state $\wbar$ with its own
$\sigma_\omega = \norm{\dot{\wbar}}_{\mathrm{rms}}/\norm{\wbar}_{\mathrm{rms}}$,
and per wavenumber shell $\kappa$ with
$\sigma(\kappa) = \norm{\dot N_\kappa}/\norm{N_\kappa}$.

\paragraph{Per-order laddering.}
Define $\sigma_m = \norm{N^{(m+1)}}_{\mathrm{rms}}/\norm{N^{(m)}}_{\mathrm{rms}}$.
Because each time derivative weights the spectrum by another factor of the
local rate,
\begin{equation}\label{eq:ladder}
\sigma_m^2
= \frac{\int \sigma(\kappa)^{2(m+1)}\,E_N(\kappa)\,d\kappa}
       {\int \sigma(\kappa)^{2m}\,E_N(\kappa)\,d\kappa},
\end{equation}
the $(m{+}1)$-th moment ratio of $\sigma(\kappa)$ under the $N$-spectrum.
Since $\sigma(\kappa)$ increases with $\kappa$, $\sigma_m$ \emph{increases
with} $m$: higher derivatives progressively sample the fast small-scale
tail. This is exactly the per-order growth escalation of the table in
\S\ref{sec:outer} ($34 \to 93$).

\paragraph{Verification.}
Analytic derivatives (spectral recursion, no finite differences); curvature
fit of the measured $\rho(\tau)$ on the first lags vs the identity
$1/\sigma$:

\begin{center}
\small
\begin{tabular}{lccccccc}
\toprule
member & $\taul(N)$ fit & $1/\sigma_N$ & $\taul\sigma_N$ &
$\taul(\omega)$ fit & $1/\sigma_\omega$ & $\taul\sigma_\omega$ &
$\sigma_0,\ \sigma_1,\ \sigma_2$\\
\midrule
Re25k   & 0.0329 & 0.0287 & $\mathbf{1.15}$ & 0.310 & 0.317 & $\mathbf{0.98}$ & 34.8,\ 60.4,\ 76.4\\
kf4     & 0.1113 & 0.1020 & $\mathbf{1.09}$ & 0.743 & 0.768 & $\mathbf{0.97}$ & 9.8,\ 20.6,\ 29.4\\
FRC-256 & ---    & 0.0670 & ---             & ---   & 0.560 & ---             & (Results CSV)\\
\bottomrule
\end{tabular}
\end{center}

The identity holds: $\taul\,\sigma = 1$ within 15\% on the $N$-process and
within 3\% on $\omega$. The residual $+10$--$15\%$ on $N$ is the expected
$O(\tau^4)$ mixture effect: $\rho$ is a superposition of per-shell
parabolas and the fast shells die first, so the pooled curve decays slightly
more slowly than the single-$\sigma$ parabola, and the curvature fit reads
slightly high. The laddering is monotone ($\sigma_0 < \sigma_1 < \sigma_2$)
in every member, as derived.

\begin{figure}[htbp]
\centering
\includegraphics[width=0.72\textwidth]{img/rho_vs_identity_Re25k.png}
\caption{FRC-Re25k: measured $\rho_N(\tau)$ and $\rho_\omega(\tau)$ against
the identity parabolas $1 - \tau^2\sigma^2/2$ with $\sigma$ computed
analytically. The parabola is not fitted to the curve; agreement at small
$\tau$ is the identity \eqref{eq:identity}.}
\label{fig:rho-re25k}
\end{figure}

\begin{figure}[htbp]
\centering
\includegraphics[width=0.72\textwidth]{img/rho_vs_identity_kf4.png}
\caption{FRC-kf4: as Fig.~\ref{fig:rho-re25k}. Wider microscale
($\taul = 0.11$ vs $0.03$): the same 27-lag window now spans only
$\sim1.2\,\taul$, so $\rho_N$ decays gently.}
\label{fig:rho-kf4}
\end{figure}

\begin{figure}[htbp]
\centering
\includegraphics[width=0.72\textwidth]{img/rho_vs_identity_FRC256.png}
\caption{FRC-256: as Fig.~\ref{fig:rho-re25k}, at $256^2$ resolution.}
\label{fig:rho-frc256}
\end{figure}

\paragraph{Two further measured facts.}
\begin{enumerate}
\item \emph{The state decorrelates an order of magnitude more slowly than
its nonlinear term:} $\sigma_N/\sigma_\omega \approx 7.5$--$11$ across
members. Over the full 27-lag span, $\rho_\omega \ge 0.96$ while $\rho_N$
falls to $0.2$--$0.6$. A time stencil applied to the $\wbar$-history
therefore sees smooth data; the difficulty of estimating $N$-derivatives
lives in the derivative ladder ($\sigma_m$ growth), not in bulk state
decorrelation.
\item \emph{$\rho_N$ plateaus at a positive value} ($\approx 0.2$ for Re25k)
instead of decaying to zero: the static component of $N$ (the
time-independent forcing plus the time mean of $-J$) never decorrelates.
This does not bias the identity --- the static part is absent from
$\norm{\dot N}$ and inflates $\norm{N}$ exactly as it inflates the measured
curvature denominator --- but integral or $1/e$ decorrelation times would be
distorted by it. $\taul$ is the clean object.
\end{enumerate}

\begin{proposition}\label{prop:microscale}
Under stationarity, $\rho(\tau) = 1 - \tau^2\sigma^2/2 + O(\tau^4)$ with
$\sigma = \norm{\dot N}_{\mathrm{rms}}/\norm{N}_{\mathrm{rms}}$ exactly;
verified to 15\% ($N$) and 3\% ($\omega$) across the measured regimes.
$\sigma_m$ grows monotonically with $m$ as the $(m{+}1)$-th moment ratio of
$\sigma(\kappa)$; verified in every member.
\end{proposition}

\section{The radius is two Taylor microscales}\label{sec:radius-micro}

Sections~\ref{sec:outer} and \ref{sec:decorr} measure two independent
timescales: the six-order series radius $\DTs$ and the one-derivative
microscale $\taul$. They are proportional:

\begin{center}
\begin{tabular}{lccc}
\toprule
member & $\DTs$ (root test) & $\taul = 1/\sigma_N$ & $C = \DTs\,\sigma_N$\\
\midrule
Re25k & 0.066 & 0.029 & 2.27\\
combo & 0.139 & 0.068 & 2.05\\
kf4   & 0.199 & 0.104 & 1.92\\
\bottomrule
\end{tabular}
\end{center}

\begin{center}
$\mathbf{C = 2.08 \pm 0.15}$, constant to 7\% across $2.5\times$ in $Re$ and
$2\times$ in $\beta$.
\end{center}

Using the directly fitted $\taul$ instead of $1/\sigma$ moves $C$ to $2.00$
(Re25k) and $1.79$ (kf4) --- the same band.

\begin{proposition}\label{prop:two-micro}
$\DTs = C\,\taul$ with $C \approx 2$ independent of $(Re,\beta)$: the
analytic closure's convergence horizon is approximately \emph{two Taylor
microscales of the nonlinear term}. All $(Re,\beta)$-dependence of $\DTs$
enters through the single measurable rate $\sigma(Re,\beta)$.
\end{proposition}

This also grounds the inner wall of \S\ref{sec:inner}: the depth-7 stencil
at Re25k starts hurting precisely when its span ($6\DT = 0.09$) exceeds
$\DTs = 2\taul$ --- i.e.\ when its oldest marks are several microscales
stale.

\section{Error propagation is benign across the envelope}\label{sec:amplification}

\paragraph{Setup.}
From \eqref{eq:defect}, $\delta = \sum_p (\DT^{\,p}/D_p)\,R_p$ with
$R_p = \sum_k a_{p,k}\,L^{p-k}\cdot\mathrm{field}_k$. The learned quantities
are the $N$-derivatives; an error $\varepsilon$ on $\mathrm{field}_k$
contributes $(\DT^{\,p}/D_p)\,|a_{p,k}|\,\varepsilon\,
\norm{L^{p-k}\,\mathrm{field}_k}$ --- terms carrying powers of $L$ are
potentially amplified.

\paragraph{Measurement 1 --- operating step $\DT = 5\times10^{-3}$}
(measured per-order errors $\varepsilon = 2.6/3.0/4.0\%$ on
$\dot N/\ddot N/\dddot N$):

\begin{center}
\begin{tabular}{lccc}
\toprule
member & $\norm{\delta}$ & dominant term & error$/\norm{\delta}$\\
\midrule
Re25k & $4.07\times10^{-5}$ & $L^0\ddot N$ ($R_3$) $= 4.07\times10^{-5}$ & 3.00\%\\
combo & $7.78\times10^{-5}$ & $L^0\ddot N = 7.77\times10^{-5}$ & 3.00\%\\
kf4   & $9.98\times10^{-6}$ & $L^0\ddot N = 9.92\times10^{-6}$ & 2.98\%\\
\bottomrule
\end{tabular}
\end{center}

The un-amplified $L^0\ddot N$ term equals $\norm{\delta}$ to three figures in
every member; the amplified terms ($L^2\dot N$, $L^1\ddot N$) are 5--7
orders of magnitude smaller.

\paragraph{Measurement 2 --- envelope edge $\DT = 1.5\times10^{-2}$}
(unit per-term errors: pure amplification geometry; Re25k, the
fastest-cascade member):

\begin{center}
\begin{tabular}{lc}
\toprule
term & contribution to $\norm{\delta}$\\
\midrule
$L^0\ddot N$ ($R_3$)  & $\mathbf{100.2\%}$\\
$L^0\dddot N$ ($R_4$) & 2.6\%\\
$L^1\dot N$ ($R_3$)   & 0.5\%\\
$L^1\ddot N$ ($R_4$)  & 0.2\%\\
$L^2\dot N$ ($R_4$)   & 0.00\%\\
\bottomrule
\end{tabular}
\end{center}

The feared $L$-amplification of the $R_4$ terms does not materialize even at
the largest trained step: the closure error tracks the $\ddot N$ error
one-to-one (correlated-field total $99.98\%$), with a $2.6\%$ additive
$\dddot N$ contribution.

\begin{proposition}\label{prop:benign}
Across the entire trained envelope $\DT \le 1.5\times10^{-2}$, the closure
error equals the relative $\ddot N$ error, un-amplified:
$\norm{\Delta\delta}/\norm{\delta} \approx \varepsilon_{\ddot N}$. Any
improvement in $\ddot N$ accuracy converts to closure-error reduction
one-to-one, and $\ddot N$ accuracy is the single quantity that sets the
rollout ceiling.
\end{proposition}

\section{Operating envelope}\label{sec:envelope}

\begin{corollary}
For a depth-$k$ stencil, operate where the stencil span stays well inside
the member's radius: $(k{-}1)\DT$ comfortably below
$\DTs(Re,\beta) = 2\taul$. Then the series converges
(Prop.~\ref{prop:radius}), the stencil resolves the $N$-derivatives
(Prop.~\ref{prop:inner}), and derivative accuracy converts to closure
accuracy one-to-one (Prop.~\ref{prop:benign}). Narrow-radius members (high
$Re$, high $\beta$) cap the pooled envelope first: at $k = 7$, Re25k caps
near $\DT \approx 10^{-2}$ while combo and kf4 remain inside through
$1.5\times10^{-2}$.
\end{corollary}

\section{Caveats}\label{sec:caveats}

\begin{itemize}
\item Finite truncation depth ($p = 3,\dots,6$) puts a $\pm20$--$30\%$ band
on $\DTs$; rankings are robust, absolute values carry the band.
\item The inner-wall crossover is directly observed only at Re25k (its wall
lies inside the measured $\DT$ range); the combo/kf4 walls lie beyond
$\DT = 1.5\times10^{-2}$ and are bounded, not located.
\item The $\taul$ identity is verified on the three members shown; decaying
(non-stationary) members require detrending before the identity applies and
are not yet measured.
\item Scope: the forced-turbulence members in the header; all statements are
for the resolved (dealiased) band.
\end{itemize}

\end{document}